\documentclass[a4paper]{article}

\usepackage[a4paper,margin=25mm]{geometry}
\usepackage[english]{babel}

\usepackage{parskip} 
\usepackage[T1]{fontenc}
\usepackage{microtype}
\usepackage{setspace}
\setstretch{1.2}


\usepackage{mathtools, amssymb}
\renewcommand{\Vec}[1]{\boldsymbol{#1}}
\DeclareMathOperator{\divergence}{div}
\DeclareMathOperator{\grad}{grad}
\DeclareMathOperator{\curl}{curl}

\usepackage{newtxtext}
\usepackage{newtxmath}


\usepackage[dvipsnames]{xcolor}
\definecolor{AUblue}{RGB}{0,61,115}
\definecolor{AUblueDark}{RGB}{0,37,70}


\usepackage{graphicx}
\graphicspath{{./figures/}}
\usepackage{tikz}
\usetikzlibrary{calc, arrows.meta}

\usepackage{enumitem}
\setlist{topsep=4pt,itemsep=2pt,parsep=0pt}

\usepackage{siunitx}
\sisetup{
  exponent-product = \cdot,
  output-decimal-marker = {,}
}

%Giles Castelles incfig
\usepackage{import}
\usepackage{xifthen}
\usepackage{pdfpages}
\usepackage{transparent}

\newcommand{\incfig}[2][1]{%
  \def\svgwidth{#1\columnwidth}
  \import{./figures/}{#2.pdf_tex}
}

\usepackage{fancyhdr}
\pagestyle{fancy}
\fancyhf{}
\renewcommand{\headrule}{\hbox to\headwidth{\color{AUblueDark}\leaders\hrule height 0.6pt\hfill}}
\setlength{\headheight}{20pt}
\fancyhead[L]{\small\AuthorName}
\fancyhead[C]{\small\ReportTitle}
\fancyhead[R]{\small\TheDate}
\usepackage{lastpage}
\fancyfoot[R]{\small\thepage{} of \pageref{LastPage}}
\fancyfoot[L]{\small\CourseName}

\newcommand{\Dept}{Department of Mechanical and Production Engineering}
\newcommand{\Uni}{Aarhus University}
\newcommand{\ReportTitle}{Assignment 1}
\newcommand{\AuthorName}{Noah Rahbek Bigum Hansen}
\newcommand{\StudentID}{202405538}
\newcommand{\CourseName}{Fluid Mechanics}
\newcommand{\CourseNumber}{290211U004}
\newcommand{\Course}{\CourseName{} --- \CourseNumber}
\newcommand{\TheDate}{October 3, 2025}
\newcommand{\Wordcount}{8,123}
\newcommand{\Logo}{figures/AU}


\usepackage{hyperref}
\hypersetup{
  colorlinks=true,
  linkcolor=AUblue,
  citecolor=AUblue,
  urlcolor=AUblue,
  pdftitle={\ReportTitle},
  pdfauthor={\AuthorName},
  pdfsubject={\Course},
  pdfcreator={LaTeX + Vimtex + pdfTeX + Libertinus},
}
\usepackage[nameinlink]{cleveref}

\usepackage{chngcntr}
\newcounter{problem}
\renewcommand{\theproblem}{\arabic{problem}}
\counterwithin{figure}{problem}
\counterwithin{table}{problem}
\numberwithin{equation}{problem}
\makeatletter
\providecommand*\theHproblem{\theproblem}
\renewcommand{\theHfigure}{problem.\theHproblem.\arabic{figure}}
\renewcommand{\theHtable}{problem.\theHproblem.\arabic{table}}
\renewcommand{\theHequation}{problem.\theHproblem.\arabic{equation}}
\makeatother

\crefname{problem}{problem}{problems}
\Crefname{problem}{Problem}{Problems}

\newenvironment{problem}%
{%
  \refstepcounter{problem}%
  \par\vspace{0pt}%
  {\large\bfseries\sffamily\color{AUblueDark}%
    Problem \theproblem%
  }%
  \par\vspace{.35\baselineskip}%
}%
{%
  \par\vspace{.6\baselineskip}%
  {\color{AUblueDark}\rule{\linewidth}{0.8pt}}%
  \vspace{.6\baselineskip}\par%
}

\newcommand{\AUheading}[1]{%
  \par\vspace{.4\baselineskip}%
  {\sffamily\bfseries\color{AUblueDark}#1\par}%
  \nobreak%
}

\newenvironment{assumptions}[1][Assumptions]{%
  \AUheading{#1}
  \begin{itemize}[leftmargin=*,nosep]
}{\end{itemize}\par\medskip}

\newenvironment{derivation}[1][Derivation]{%
  \AUheading{#1}
  \begingroup\allowdisplaybreaks
  \setlength{\abovedisplayskip}{6pt}%
  \setlength{\belowdisplayskip}{6pt}%
  \setlength{\abovedisplayshortskip}{4pt}%
  \setlength{\belowdisplayshortskip}{4pt}%
}{\par\endgroup\medskip}

\newcommand{\finalresult}[1]{%
  \vspace{.2\baselineskip}%
  \begin{equation}
    {\color{AUblueDark}\boxed{\displaystyle #1}}
    \tag*{\textsc{\textcolor{AUblueDark}{Result}}}
  \end{equation}
  \newpage
}

\usepackage{environ}
\newlength{\pbfigsep}
\setlength{\pbfigsep}{8mm}

\NewEnviron{problemwithimage}[2][0.38\linewidth]{%
  \begin{problem}%
    \noindent
    \begin{minipage}[t]{\dimexpr\linewidth-#1-\pbfigsep\relax}
      \vspace{-0.2\baselineskip}\BODY
    \end{minipage}\hfill
    \begin{minipage}[t]{#1}
      \vspace{-\baselineskip}\centering
      \includegraphics[width=\linewidth]{#2}%
    \end{minipage}%
  \end{problem}%
}

\begin{document}
\begin{titlepage}
\thispagestyle{empty}
\onehalfspacing

\begin{tikzpicture}[remember picture,overlay]
  \fill[AUblue] ([xshift=-8mm]current page.west |- current page.north)
       rectangle ([xshift=7mm]current page.west |- current page.south);
\end{tikzpicture}

\vspace*{8mm}
{\large\bfseries \Dept}\\[2pt]
{\normalsize \Uni}
\hfill

\vspace{28mm}
{\fontsize{30pt}{32pt}\selectfont\bfseries \ReportTitle\par}
\vspace{12mm}
{\color{AUblueDark}\rule{\linewidth}{1.5pt}}

\vspace{10mm}
\begin{tabular}{@{}ll}
\textbf{Author:}     & \AuthorName \\
\textbf{Student ID:} & \StudentID \\
\textbf{Course:}     & \Course \\
\textbf{Date:}       & \TheDate \\
\end{tabular}

\begin{tikzpicture}[remember picture,overlay]
  \node[anchor=south east, inner sep=0] at
    ($ (current page.south east) + (-20mm, 20mm) $)
    {\includegraphics[width=42mm]{\Logo}};
\end{tikzpicture}

\vfill
{\small Word count: \Wordcount}
\end{titlepage}

\setcounter{page}{1}

\begin{problem}
Water flows steadily through a fire hose and nozzle.
The hose is \qty{35}{mm} in diameter and the nozzle tip is \qty{25}{mm} in diameter; water gauge pressure in the hose is \qty{510}{kPa}, and the stream leaving the nozzle is uniform.
The exit speed and pressure are \qty{32}{m/s} and atmospheric, respectively.

Find the force transmitted by the coupling between the nozzle and hose.
Indicate whether the coupling is in tension or compression.
State the assumptions made and sketch the situation.
\end{problem}
\begin{figure}[ht]
    \centering
    \incfig{p1_1}
    \caption{Sketch of the coupling.}
    \label{fig:p1_1}
\end{figure}

\begin{assumptions}
  \item Steady flow
  \item Uniform flow at interfaces between coupling, hose, and nozzle
  \item Incompressible and frictionless liquid with negligible mass
\end{assumptions}

\begin{derivation}
The momentum equation for an inertial control volume (as is the case for the coupling) is
\[ 
\Vec{F} = \Vec{F}_B + \Vec{F}_S + \frac{\partial }{\partial t} \int_{\mathrm{CV}} \Vec{V} \rho \, \mathrm{d}V + \int_{\mathrm{CS}} \Vec{V}\rho \Vec{V} \cdot \mathrm{d}\Vec{A}
.\]
In this case the flow is steady, hence $\frac{\partial }{\partial t} = 0$, hence the $\Vec{F}_B = 0$. Reducing the vector equation to a scalar equation in the $x$-direction (defined on \Cref{fig:p1_1}) we obtain
\[ 
F_x = F_{S_x} = \int_{CS} u \rho \Vec{V} \cdot \mathrm{d}\Vec{A}
.\]
As the flow is assumed uniform at both the inlet and outlet the integral further reduces to a sum as
\begin{equation}\label{eq:momste1d}
  F_x = F_{S_x} = \sum_{\mathrm{CS}} u \rho \Vec{V} \cdot \Vec{A}
\end{equation}
The forces in the $x$-direction on the control volume are the pressure at the interface between the hose and the coupling ($+p_h A_h$), the pressure at the nozzle ($-p_n A_n$), and the reaction on the coupling $R_x$. From \Cref{eq:momste1d} we get
\begin{equation} \label{eq:momste1dsub}
  R_x + p_h A_h - p_n A_n = \rho Q \left( V_n - V_h \right)
\end{equation}
As we are using gauge pressures $p_n = \qty{0}{kPa} \implies p_n A_n = \qty{0}{N}$, therefore \Cref{eq:momste1dsub} reduces to
\begin{equation} \label{eq:momste1dsubred}
  R_x + p_h A_h = \rho Q \left( V_n - V_h \right)
\end{equation}
The continuity equation prescribes that the flow rate $Q = V A = \mathrm{const.}$ Hence,
\begin{subequations} \label{eq:cont}
  \begin{align}
  Q &= A_n V_n\\
  V_h &= \frac{Q}{A_h} = \frac{A_n V_n}{A_h}
\end{align}
\end{subequations}
Substituting \Cref{eq:cont} into \Cref{eq:momste1dsubred} and solving for $R_x$ we obtain
\begin{align*}
  R_x + p_h A_h &= \rho A_n V_n \left( V_n - \frac{A_n V_n}{A_h} \right) \\
  R_x &= \rho A_n V_n^2 \left( 1 - \frac{A_n}{A_h} \right) - p_h A_h \\
  R_x &= \qty{1000}{\frac{kg}{m^3}} \cdot \left( \pi \cdot \left( \qty{12,5}{mm}  \right)^2 \right) \cdot \left( \qty{32}{\frac{m}{s}}  \right)^2 \left( 1 - \frac{ \left( \qty{12,5}{mm}  \right)^2 }{\left( \qty{17,5}{mm}  \right)^2} \right) - \qty{510}{kPa} \cdot \pi \cdot \left( \qty{17,5}{mm} \right)^2 \\
  R_x &= \qty{-244,48}{N} 
.\end{align*}
The negative sign (with $+x$ being to the right as shown on \Cref{fig:p1_1}) means the coupling reaction force points to the left, and the joint will therefore be in a tension of about \qty{245}{N}.
\end{derivation}

\finalresult{R_c = \qty{245}{N} \hspace{0.5em} \left( \text{Tension} \right)}


\begin{problemwithimage}[0.35\linewidth]{./figures/p2_1.png}
  Determine the gauge pressure in \unit{kPa} at point $a$, if liquid $A$ has $\mathrm{SG} = \num{1,20} $ and liquid $B$ has $\mathrm{SG} = \num{0,75}$.
  The liquid surrounding point $a$ is water, and the tank on the left is open to the atmosphere.
\end{problemwithimage}

\begin{problem}
 Determine the stream function $\psi$ that will yield the velocity field.
\[ 
\Vec{V} = 2y \left( 2x + 1 \right) \hat{i} + \left( x \left( x + 1 \right) - 2y^2 \right) \hat{j} \, V
.\] 
\end{problem}


\begin{problemwithimage}[0.5\linewidth]{./figures/p4_1.png}
Water flows steadily through the nozzle, shown in the figure below, discharging to atmosphere. Calculate the horizontal component of force in the flanged joint.
Indicate whether the joint is in tension or compression.
Use density of $\mathrm{water} = \qty{1000}{\frac{kg}{m^3}}$.
\end{problemwithimage}


\begin{problem}
The velocity components for an incompressible steady flow field are given below
\[ 
u = A \left( x^2 + z^2 \right) \qquad v = B \left( xy + yz \right)
.\]
Determine the $z$-component of velocity for steady and unsteady flow.
\end{problem}


\begin{problemwithimage}[0.35\linewidth]{./figures/p6_1.png}
The water flow rate through the vertical bend shown in the figure below is $\qty{2,83}{\frac{m^3}{s}}$.
Calculate the magnitude, direction, and location of the resultant force of the water pipe bend.
Use density of $\mathrm{water}= \qty{999}{\frac{kg}{m^3}}$.
\end{problemwithimage}



\begin{problem}
A viscous liquid is sheared between two parallel disks of radius $R$, one of which rotates while the other is fixed.
The velocity field is purely tangential, and the velocity varies linearly with $z$ from $V_{\phi} = 0$ at $z = 0$ (the fixed disk) to the velocity of the rotating disk at its surface ($z = h$).

Derive an expression for the velocity field between the disks.
\end{problem}


\end{document}

